\section*{Bab \ref{ch:b2:magmatic-differentiation} --- Diferensiasi Magma}

\begin{solution}{exo:b2:magmatic-differentiation:1}
(a) Likuidus An$_{50}$ adalah \qty{1450}{\celsius}. Kristal pertama terletak pada
solidus pada suhu itu, di antara An$_{66.7}$ (1372) dan An$_{83.3}$ (1465):
$66.7 + 16.6\times(1450 - 1372)/(1465 - 1372) = 66.7 + 13.9 = \qty{80.6}{\percent}$
anortit. (b) Pada \qty{1394}{\celsius}, cairan berada pada likuidus, yaitu
An$_{33.3}$; kristal berada pada solidus, $66.7 + 16.6\times22/93 =
\qty{70.6}{\percent}$. Menurut aturan tuas, fraksi yang telah mengkristal adalah $(50 -
33.3)/(70.6 - 33.3) = 16.7/37.3 = 0.448$. (c) Cairan terakhir lenyap ketika
kristal mencapai An$_{50}$, pada suhu solidus \qty{1287}{\celsius}; likuidus
di sana terletak di antara An$_{11.1}$ (1265) dan An$_{20}$ (1334): $11.1 +
8.9\times22/69 = \qty{13.9}{\percent}$ anortit.
\end{solution}

\begin{solution}{exo:b2:magmatic-differentiation:2}
$\bar D_{\ce{Sr}} = 0.15\times0.000063 + 0.35\times0.157 + 0.50\times2.7 = 0.00001 +
0.0550 + 1.35 = 1.40$; $\bar D_{\ce{Zr}} = 0.0001 + 0.0683 + 0.0005 = 0.069$;
$\bar D_{\ce{Rb}} = 0.00001 + 0.0012 + 0.0125 = 0.014$. Hanya stronsium, yang ditampung oleh
plagioklas, yang kompatibel ($\bar D > 1$); zirkonium dan rubidium inkompatibel,
dan rubidium lebih inkompatibel lagi.
\end{solution}

\begin{solution}{exo:b2:magmatic-differentiation:3}
(a) $C_L/C_0 = 0.5^{D-1}$: Sr $0.5^{0.40} = 0.758$ (turun); Zr $0.5^{-0.931} =
1.907$; Rb $0.5^{-0.986} = 1.981$, mendekati limit 2. (b) Zr/Rb berubah dengan faktor
$0.5^{0.069 - 0.014} = 0.5^{0.055} = 0.963$, kurang dari \qty{4}{\percent}, sementara setiap
unsur hampir menjadi dua kali lipat: seperti dinyatakan korolari, rasio dua \omterm{def:b2:elements-in-the-earth:partition}{unsur
inkompatibel} mempertahankan nilai \omterm{def:b2:magma-emplacement:magma}{magma} induk, sehingga memberi tahu tentang sumber, bukan
tentang kristalisasi.
\end{solution}

\begin{solution}{exo:b2:magmatic-differentiation:4}
(a) \ce{SiO2} $+5.03$, \ce{MgO} $-1.95$, \ce{K2O} $+0.82$ (persen); kalium oksida
dikalikan $1.67/0.85 = 1.96$. (b) Untuk unsur yang sepenuhnya inkompatibel,
$C_L/C_0 = 1/F$, sehingga $F = 0.85/1.67 = 0.509$: separuh \omterm{def:b2:magma-emplacement:magma}{magma} harus sudah
mengkristal. (c) Penyingkiran \qty{12.5}{g} per \qty{100}{g} menyisakan $F \approx 0.88$ dan
hanya akan menaikkan kalium oksida dengan faktor sekitar $1/0.88 = 1.14$: kristalisasi saja tidak dapat
menjelaskan kenaikan itu. Kalium oksida datang dari luar, dari xenolit granitik yang kaya akan
kalium oksida (\qty{3.6}{\percent}), yang diasimilasi oleh \omterm{def:b2:magma-emplacement:magma}{magma}.
\end{solution}

\begin{solution}{exo:b2:magmatic-differentiation:5}
(a) $F = 0.8$, $D = 5$. Fraksional: $0.8^{4} = 0.410$; setimbang:
$1/(5 + 0.8\times(1 - 5)) = 1/1.8 = 0.556$. (b) $F = 0.1$, $D = 0.01$. Fraksional:
$0.1^{-0.99} = 9.77$; setimbang: $1/(0.01 + 0.1\times0.99) = 1/0.109 = 9.17$.
(c) Ketika $F \to 0$, hukum Rayleigh menuju $0$ untuk $D > 1$ dan menuju tak hingga untuk
$D < 1$, sedangkan hukum setimbang menuju $1/D$ (cairan terakhir setimbang
dengan kristal yang berkomposisi total). Hanya \omterm{def:b2:magmatic-differentiation:crystallisation}{kristalisasi fraksional} yang dapat melucuti suatu
\omterm{def:b2:elements-in-the-earth:partition}{unsur kompatibel} dari cairan sepenuhnya.
\end{solution}

\begin{solution}{exo:b2:magmatic-differentiation:6}
(a) Satu satuan rumus mengandung 1.6 \ce{MgO}, 0.4 \ce{FeO}, dan satu \ce{SiO2}: $1.6\times
40.30 = 64.48$, $0.4\times71.84 = 28.74$, 60.08, jumlahnya \qty{153.30}{g}. Persentase
massa: \ce{MgO} \qty{42.06}{\percent}, \ce{FeO} \qty{18.75}{\percent},
\ce{SiO2} \qty{39.19}{\percent}. (b) Aturan tuas pada \ce{MgO}: $x = (10.0 - 8.0)/(42.06
- 8.0) = 2.0/34.06 = 0.0587$, sekitar \qty{6}{\percent} olivin. (c) Silika
cairan: $(49.0 - 0.0587\times39.19)/(1 - 0.0587) = 46.70/0.9413 =
\qty{49.61}{\percent}$. Penyingkiran olivin, yang lebih miskin silika daripada basal, menaikkan
silika cairan, tetapi hanya perlahan.
\end{solution}

\begin{solution}{exo:b2:magmatic-differentiation:7}
(a) $F(0) = 0.768$; $F(1000) = 0.1091 - 0.5906 + 0.7678 = 0.286$; $F(2000) = 0.4364
- 1.1813 + 0.7678 = 0.023$. (b) Diskriminannya $(5.9064\times10^{-4})^2 -
4\times1.091\times10^{-7}\times0.7678 = 3.4885\times10^{-7} - 3.3507\times10^{-7} =
1.379\times10^{-8}$, dengan akar kuadrat $1.174\times10^{-4}$; akar yang lebih kecil adalah
$H = (5.9064 - 1.174)\times10^{-4}/(2.182\times10^{-7}) = \qty{2169}{m}$, yaitu
Horizon Sandwich. (c) $1 - 0.768 = 0.232$: hampir seperempat \omterm{def:b2:magma-emplacement:magma}{magma} telah
mengkristal sebelum dasar Zona Bawah yang tersingkap terbentuk, di Zona Tersembunyi
di bawahnya dan di Seri Tepi di sepanjang dinding dan atap.
\end{solution}

\begin{solution}{exo:b2:magmatic-differentiation:8}
\omterm{def:b2:magmatic-differentiation:cumulate}{Kumulat} mengandung $D\,C_L = D\,C_0F^{D-1}$, sehingga di antara dua tingkat
$\ln(\mathrm{Ni}_2/\mathrm{Ni}_1) = (D - 1)\ln(F_2/F_1)$. (a) $D - 1 = \ln(54/177)/
\ln(0.29/0.70) = (-1.187)/(-0.881) = 1.35$, sehingga $D \approx 2.35$ (regresi atas
66 \omterm{def:b2:magmatic-differentiation:cumulate}{kumulat} dengan $F \ge 0.23$ memberikan 2.38). (b) $D - 1 = \ln(6.2/54)/\ln(0.17/0.29) = (-2.164)/(-0.534)
= 4.05$, sehingga $D \approx 5.05$. (c) $D$ total berubah ketika himpunan mineral yang
mengkristal berubah (mineral baru yang menyerap nikel, atau mineral yang berhenti mengkristal), dan
koefisien mineral--lelehan sendiri naik ketika cairan mendingin dan berubah
komposisi, seperti yang terukur pada olivin; cairan yang terperangkap dalam \omterm{def:b2:magmatic-differentiation:cumulate}{kumulat} mengaburkan
taksiran itu.
\end{solution}

\begin{solution}{exo:b2:magmatic-differentiation:9}
(a) $100 + 25.4 - 2.9 - 9.6 = \qty{112.9}{g}$. (b) Gram \ce{SiO2}: $55.53 +
25.4\times0.7248 - 2.9\times0.3919 - 9.6\times0.5054 = 55.53 + 18.41 - 1.14 - 4.85 =
\qty{67.95}{g}$, sehingga $67.95/112.9 = \qty{60.19}{\percent}$ dibandingkan dengan 60.16. \ce{MgO}:
$4.66 + 0.20 - 1.22 = \qty{3.64}{g}$, \qty{3.22}{\percent} dibandingkan dengan 3.19. \ce{K2O}:
$0.88 + 0.92 = \qty{1.80}{g}$, \qty{1.60}{\percent} dibandingkan dengan 1.75. (c) Silika dan magnesia, dua dari tiga
pengontrol, cocok sampai selisih 0.03; kalium oksida, yang tidak dipakai sebagai pengontrol, lebih rendah 0.15, hampir
sepersepuluh nilainya: bahan yang terasimilasi mengandung lebih banyak kalium oksida daripada xenolit
rata-rata, atau lelehannya yang kaya kalium oksida lebih banyak terambil.
\end{solution}

\begin{solution}{exo:b2:magmatic-differentiation:10}
(a) $z = (0.2 + 1.5 - 1)/(0.2 - 1) = 0.7/(-0.8) = -0.875$; $F^{-z} = 0.7^{0.875} =
0.7319$; $C^*/C_0 = 0.2\times0.5/(0.2 - 1 + 1.5) = 0.1/0.7 = 0.1429$, sehingga $C/C_0 =
0.7319 + 0.1429\times(1 - 0.7319) = 0.7319 + 0.0383 = 0.770$. (b) Rayleigh:
$0.7^{0.5} = 0.837$. \omterm{def:b2:magma-emplacement:magma}{Magma} AFC mengandung lebih sedikit stronsium karena dua alasan: untuk kehilangan
\qty{30}{\percent} massanya sambil memperoleh $r$ kali massa yang dikristalkannya, \omterm{def:b2:magma-emplacement:magma}{magma} itu harus
mengkristalkan $(1 - 0.7)/(1 - 0.2) = 0.375$, bukan 0.30, dan stronsium yang diperolehnya
berasal dari batuan yang kadarnya separuh. (c) $(0.7319\times0.7035 + 0.0383\times0.7200)/0.7702
= (0.5149 + 0.0276)/0.7702 = 0.70432$: rasio itu naik 0.0008, padahal
kristalisasi saja akan membiarkannya tetap 0.7035.
\end{solution}

\begin{solution}{exo:b2:magmatic-differentiation:11}
(a) $1.1\times(1100 - 200) + 250 = 990 + 250 = \qty{1240}{J}$ per gram. (b)
$25.4\times1240 = \qty{31496}{J}$, mendekati \qty{31350}{J} Wilcox. (c)
$31496/370 = \qty{85.1}{g}$ plagioklas. (d) Neraca massa hanya menyingkirkan $2.9 +
9.6 = \qty{12.5}{g}$ kristal, yang kalor latennya (sekitar \qty{4600}{J}) menutupi kurang
dari seperenam kebutuhan; $r = 25.4/12.5 = 2.03$, batuan yang terasimilasi dua kali lebih banyak daripada
kristal yang terbentuk, padahal neraca kalor hanya mengizinkan $r$ jauh di bawah 1. \omterm{def:b2:magma-emplacement:magma}{Magma} yang meletus
tidak mungkin memasok kalor itu sendiri: massa \omterm{def:b2:magma-emplacement:magma}{magma} yang lebih besar pasti telah mengkristal
dan mendingin di kedalaman, tanpa terlihat.
\end{solution}

\begin{solution}{exo:b2:magmatic-differentiation:12}
(a) Tahap 1, $F$ dari 1 ke 0.8 dengan $D = 5$: faktor $0.8^{4} = 0.4096$. Tahap 2,
dari 0.8 ke 0.5 dengan $D = 1.5$: hukum itu berlaku pada cairan yang tersisa, sehingga
faktornya $(0.5/0.8)^{0.5} = 0.7906$. Nikel: $250\times0.4096\times0.7906 =
\qty{81.0}{ppm}$. (b) $250\times0.5^{4} = \qty{15.6}{ppm}$, lima kali terlalu rendah. (c)
Hukum Rayleigh mengandaikan $D$ yang tetap; ketika himpunan mineral yang mengkristal berubah,
$D$ melompat, dan hukum itu harus dimulai lagi dari komposisi cairan yang dicapai pada
perubahan itu, dengan $F$ diukur relatif terhadap cairan tersebut.
\end{solution}

\begin{solution}{pb:b2:magmatic-differentiation:1}
\noindent\textbf{Bagian I.}

\textbf{1.} $\bar D_{\ce{Ni}} = 0.20\times5.0 + 0.50\times0 + 0.30\times2.0 = 1.60$.

\textbf{2.} $\bar D_{\ce{Zr}} = 0.20\times0.0007 + 0.50\times0.001 + 0.30\times0.195 =
0.00014 + 0.0005 + 0.0585 = 0.059$.

\textbf{3.} Zona atas: $\bar D_{\ce{Ni}} = 0.45\times2.0 = 0.90$; $\bar D_{\ce{Zr}} =
0.55\times0.001 + 0.45\times0.195 = 0.088$.

\textbf{4.} Nikel kompatibel di zona bawah (1.60) dan menjadi sedikit
inkompatibel (0.90) ketika olivin menghilang; zirkonium inkompatibel di kedua zona.

\textbf{5.} Klinopiroksen: mineral ini menyumbang $0.0585$ dari $0.059$ di zona bawah, hampir
seluruhnya, karena hanya mineral inilah yang koefisiennya jauh dari nol.

\textbf{6.} $(C_{\ce{Ni}}/C_{\ce{Zr}})$ berubah dengan faktor $F^{D_{\ce{Ni}} - D_{\ce{Zr}}}$;
untuk setiap pengurangan separuh, $0.5^{1.60 - 0.059} = 0.5^{1.541} = 0.344$.

\noindent\textbf{Bagian II.}

\textbf{7.} $1.60\times150 = \qty{240}{ppm}$.

\textbf{8.} Cairan saat itu mengandung $190/1.60 = \qty{118.75}{ppm}$; $F = (118.75/150)^{1/0.6} =
0.7917^{1.667} = 0.68$.

\textbf{9.} $80\times0.68^{-0.941} = 80\times1.438 = \qty{115}{ppm}$.

\textbf{10.} Fraksional: $150\times0.68^{0.6} = 150\times0.793 = \qty{119.0}{ppm}$;
setimbang: $150/(1.6 + 0.68\times(1 - 1.6)) = 150\times0.839 = \qty{125.8}{ppm}$.
Selisihnya hanya \qty{6}{\percent} setelah sepertiga \omterm{def:b2:magma-emplacement:magma}{magma}
mengkristal: untuk unsur yang cukup kompatibel, kedua hukum itu tetap berdekatan pada awalnya.

\textbf{11.} $150\times0.55^{0.6} = \qty{104.8}{ppm}$.

\textbf{12.} Tepat di bawah: $1.60\times104.8 = \qty{167.7}{ppm}$; tepat di atas:
$0.90\times104.8 = \qty{94.3}{ppm}$. Profil nikel \omterm{def:b2:magmatic-differentiation:cumulate}{kumulat} turun
tiba-tiba sebesar \qty{44}{\percent} pada tingkat tempat olivin menghilang dari
lapisan-lapisan: sebuah undakan, bukan kurva yang mulus.

\textbf{13.} Cairan: $104.8\times(0.40/0.55)^{0.90 - 1} = 104.8\times1.032 =
\qty{108.2}{ppm}$; \omterm{def:b2:magmatic-differentiation:cumulate}{kumulat} $0.90\times108.2 = \qty{97.4}{ppm}$. Dengan $\bar D < 1$,
nikel kini naik perlahan dalam cairan dan dalam \omterm{def:b2:magmatic-differentiation:cumulate}{kumulat}: profilnya, setelah
undakan itu, naik landai.

\noindent\textbf{Bagian III.}

\textbf{14.} $x_{\mathrm{ol}}(m_{\mathrm{ol},j} - d_j) + x_{\mathrm{pl}}
(m_{\mathrm{pl},j} - d_j) + x_{\mathrm{cpx}}(m_{\mathrm{cpx},j} - d_j) = 100\,(p_j -
d_j)$.

\textbf{15.} \ce{MgO}: $37.16\,x_{\mathrm{ol}} - 4.9\,x_{\mathrm{pl}} +
11.1\,x_{\mathrm{cpx}} = 410$; \ce{CaO}: $-7.3\,x_{\mathrm{ol}} +
7.06\,x_{\mathrm{pl}} + 13.2\,x_{\mathrm{cpx}} = 420$; \ce{Al2O3}:
$-14.1\,x_{\mathrm{ol}} + 17.6\,x_{\mathrm{pl}} - 11.1\,x_{\mathrm{cpx}} = 190$.

\textbf{16.} Penyelesaian (dengan eliminasi atau aturan Cramer): $x_{\mathrm{ol}} = 9.03$,
$x_{\mathrm{pl}} = 30.84$, $x_{\mathrm{cpx}} = 20.32$ gram per \qty{100}{g}.

\textbf{17.} Kristal $9.03 + 30.84 + 20.32 = \qty{60.19}{g}$; cairan
\qty{39.81}{g}: $F = 0.398$.

\textbf{18.} $(4950 - 9.03\times39.19 - 30.84\times50.54 - 20.32\times51.5)/39.81 =
(4950 - 353.9 - 1558.7 - 1046.5)/39.81 = 1990.9/39.81 = \qty{50.0}{\percent}$, sama seperti pada
\omterm{def:b2:magma-emplacement:bodies}{retas}: oksida yang tidak dipakai sebagai pengontrol ikut tereproduksi.

\textbf{19.} Kristal yang tersingkir terdiri atas $9.03/60.19 = \qty{15}{\percent}$ olivin,
\qty{51}{\percent} plagioklas, dan \qty{34}{\percent} klinopiroksen, di antara
komposisi modal kedua zona. Kristal itu cocok dengan himpunan bawah yang mengkristal dari $F = 1$ sampai
0.55 ($0.45\times20 = 9.0$, $0.45\times50 = 22.5$, $0.45\times30 = 13.5$) yang disusul oleh
himpunan atas dari 0.55 sampai 0.40 ($0.15\times55 = 8.25$, $0.15\times45 = 6.75$): 9.0,
30.75, dan 20.25 gram.

\textbf{20.} $0.30/0.75 = 0.40$.

\noindent\textbf{Bagian IV.}

\textbf{21.} $0.55^{0.059 - 1} = 0.55^{-0.941} = 1.755$; $80\times1.755 =
\qty{140.4}{ppm}$.

\textbf{22.} $188/140.4 = 1.339 = (F/0.55)^{0.088 - 1}$, sehingga $F = 0.55\times1.339^{-1/
0.912} = 0.55\times0.726 = 0.399$.

\textbf{23.} $(188/80)^{1/(0.07 - 1)} = 2.35^{-1.075} = 0.399$. Untuk \omterm{def:b2:elements-in-the-earth:partition}{unsur
inkompatibel}, $C_L/C_0 \approx 1/F$ berapa pun $D$ yang kecil itu: jawabannya hampir tidak bergantung pada
rincian $D$, sehingga \omterm{def:b2:elements-in-the-earth:partition}{unsur inkompatibel} menjadi pengukur $F$ yang terbaik.

\textbf{24.} Plagioklas, yang tidak mengandung besi, merupakan separuh kristal, dan
piroksen serta olivin bersama-sama mengandung lebih sedikit besi daripada magnesia: cairan
menjadi lebih kaya akan besi dan lebih miskin magnesia tanpa bertambah silika. Pengayaan
besi pada silika yang hampir tetap ini adalah tanda seri toleiitik pada
bab berikutnya.

\textbf{25.} Kristal kumulus murni akan mengandung $0.088\times188 = \qty{16.5}{ppm}$. Dengan
fraksi massa $t$ cairan terperangkap, $16.5(1 - t) + 188\,t = 47$, sehingga $t = (47 -
16.5)/(188 - 16.5) = 30.5/171.5 = 0.18$: sekitar \qty{18}{\percent} cairan terperangkap.

\textbf{26.} Ketiga taksiran, yaitu unsur utama (0.398), kalium oksida (0.40), dan
zirkonium (0.399), sepakat: ketika lapisan teratas mengkristal,
\textbf{\qty[text-series-to-math=true]{40}{\percent}} \omterm{def:b2:magma-emplacement:magma}{magma} asal masih
cair.
\end{solution}
